\section{Engine \& Drivetrain Systems}

This section covers engine torque curves, transmission ratios, clutch engagement, and the calculation of drive forces applied to the wheels.

\subsection{Torque Curve (Parabolic)}

\begin{equation}
  T(rpm) = \max\left(0, 1 - 2.5 \left(\frac{rpm - rpm_{\text{peak}}}{rpm_{\text{range}}}\right)^2\right)
  \label{eq:torque-curve}
\end{equation}

\textbf{Physical Meaning}: The normalized torque curve is a parabola with peak output (1.0) at a target RPM. The quadratic falloff allows gentle power delivery at idle and redline, making the engine driveable across all gears.

\textbf{Source Code}: \coderef{physics.js}{301--308}

\begin{lstlisting}[caption={Parabolic torque curve}]
const rpmRange = redlineRpm - idleRpm;
const distanceFromPeak = (rpm - TORQUE_PEAK_RPM) / rpmRange;
return Math.max(0, 1.0 - 2.5 * distanceFromPeak * distanceFromPeak);
\end{lstlisting}

\textbf{Parameters}:
\begin{itemize}
  \item $rpm_{\text{range}} = rpm_{\text{redline}} - rpm_{\text{idle}}$
  \item $rpm_{\text{peak}} = \mathrm{TORQUE\_PEAK\_RPM}$ (default: 5500 RPM)
  \item Coefficient 2.5 controls falloff steepness
\end{itemize}

\subsection{Clutch Engagement Curve}

\begin{equation}
  \text{engagement}(x) = \begin{cases}
    0 & \text{if } x < b_{\text{point}} \\
    \left(\frac{x - b_{\text{point}}}{b_{\text{range}}}\right)^n & \text{if } b_{\text{point}} \leq x < b_{\text{point}} + b_{\text{range}} \\
    1 & \text{if } x \geq b_{\text{point}} + b_{\text{range}}
  \end{cases}
  \label{eq:clutch-engagement}
\end{equation}

\textbf{Physical Meaning}: Pedal position is mapped to engagement via a piecewise function with a power-law ramp in the "bite zone." This creates realistic clutch feel with a narrow, responsive engagement point.

\textbf{Source Code}: \coderef{physics.js}{323--334}

\textbf{Parameters}:
\begin{itemize}
  \item $b_{\text{point}}$: Start of bite zone (default: 0.15)
  \item $b_{\text{range}}$: Width of bite zone (default: 0.3)
  \item $n$: Power curve exponent (default: 1.5)
\end{itemize}

\subsection{Engine Torque and Drive Force}

\begin{equation}
  \tau_{\text{engine}} = T_{\text{peak}} \cdot T(rpm) \cdot \theta_{\text{throttle}}
  \label{eq:engine-torque}
\end{equation}

\begin{equation}
  F_{\text{drive}} = \frac{\tau_{\text{engine}} \cdot |r_{\text{gear}}| \cdot r_{\text{final}} \cdot e_{\text{clutch}}}{R_{\text{wheel}}}
  \label{eq:drive-force}
\end{equation}

\textbf{Physical Meaning}: Engine torque is scaled by the normalized curve, throttle input, and transmission reduction ratios. Drive force is computed by converting wheel torque to force via the wheel radius.

\textbf{Source Code}: \coderef{physics.js}{549--575}

\textbf{Parameters}:
\begin{itemize}
  \item $T_{\text{peak}}$: Peak engine torque in N·m (default: 400 N·m)
  \item $r_{\text{gear}}$: Current gear ratio (e.g., 3.6 for 1st gear)
  \item $r_{\text{final}}$: Final drive ratio (default: 3.5)
  \item $e_{\text{clutch}}$: Clutch engagement [0, 1]
  \item $R_{\text{wheel}}$: Wheel radius (default: 0.35 m)
\end{itemize}

\subsection{Idle Creep}

\begin{equation}
  F_{\text{creep}} = F_{\text{idle}} \cdot e_{\text{clutch}} \quad \text{(if } \theta < 0.05 \text{ and } v < 3 \text{ m/s)}
  \label{eq:idle-creep}
\end{equation}

\textbf{Physical Meaning}: At idle with the clutch engaged and no throttle input, a small creep force propels the car forward at low speeds, mimicking real vehicle behavior.

\textbf{Source Code}: \coderef{physics.js}{569--574}

\textbf{Parameters}:
\begin{itemize}
  \item $F_{\text{idle}}$: Idle creep force (default: 50 N)
  \item Threshold: throttle $< 5\%$, speed $< 3$ m/s
\end{itemize}

\newpage
